% Deliverable 3 fresh-start draft. Does not modify or replace any D2 report file.
% This section should only be integrated into report/main.tex after the D3
% experiments have been re-run and the tables reviewed for provenance.
\section{Phase 5: Comparative Analysis and Failure Diagnostics}
\label{sec:d3-phase5}

\subsection{Research Questions and Method}
The purpose of this phase is to evaluate the search behaviour of learned
heuristics relative to Manhattan A$^\ast$ on the same held-out mazes. It
separates prediction quality from search performance and records the exact
model backend, trained-weight artefact, split, and experimental settings.

\subsection{Paired Search Comparisons}
% After running scripts/failure_analysis.py and scripts/make_phase5_report_assets.py:
% \begin{table}[ht]
\centering
\small
\begin{tabular}{lrrr}
\toprule
Heuristic & Mazes & Mean $\Delta$ nodes & Node ratio \\
\midrule
neural\_torch\_raw & 282 & 317.784 & 3.851 \\
\bottomrule
\end{tabular}
\caption{Paired comparisons against Manhattan A*, calculated on matched test mazes.}
\label{tab:d3_paired}
\end{table}

% ADD: discussion of paired differences, runtime, path optimality, CI and p-values.

\subsection{Feature Ablations}
% After running scripts/run_feature_ablation.py and regenerating tables:
% \begin{table}[ht]
\centering
\small
\begin{tabular}{lrrr}
\toprule
Variant & Features & MAE & $R^2$ \\
\midrule
full\_29 & 29 & 0.077 & 0.850 \\
without\_wall\_window & 4 & 0.077 & 0.844 \\
goal\_offsets\_only & 2 & 0.079 & 0.845 \\
without\_goal\_offsets & 27 & 0.298 & 0.046 \\
without\_coordinates & 27 & 0.079 & 0.849 \\
\bottomrule
\end{tabular}
\caption{Fresh feature-ablation results from the original D2 data split.}
\label{tab:d3_features}
\end{table}

% ADD: compare geometry-only and local-wall-context features with equal budgets.

\subsection{Heuristic Consistency and Error Patterns}
% ADD: generated figures and measurements from the trained model under study.
% Consider adjacent-cell consistency, overestimation, maze density and dead ends.

\subsection{Checkpoint Provenance and Reproducibility}
% ADD: results/analysis/d3/torch_checkpoint_audit.json, exact SHA256 and metadata.
% Distinguish the original D2 multi-seed training from the separately benchmarked
% scikit-learn fallback. Do not describe trained checkpoints as smoke artefacts.

\subsection{Interpretation and Limitations}
% ADD ONLY AFTER VALIDATING NEW EXPERIMENTS.
% A poor search result is still scientifically useful if the analysis explains it.
